\makeabstract
    {
    Community Perspectives on Civic Engagement: Findings from Amherst{\textquoteright}s Civic Academy
    }
    {
    Litzy Casta\~neda\textsuperscript{1}, Kiarra Barnes\textsuperscript{1}, Sarah Gillespie\textsuperscript{1}
    }
    {
    \textsuperscript{1} Department of Psychology, Amherst College, Amherst, MA
    }
    {
    Civic engagement can take many forms, including voting, volunteering, attending town meetings, and working with others to address local concerns. This study looks at how Amherst residents understand their role in the community and whether participating in the Town of Amherst{\textquoteright}s Civic Academy, an 8-week, in-person introduction to town government, affects how connected they feel to Amherst, how much influence they believe they have, and how they view local government.
    }
    {
    We interviewed 28 Amherst residents: 10 Civic Academy participants, 8 residents who applied but were waitlisted, and 10 other community members with similar backgrounds. Interviews lasted between 45 minutes and two hours. Prompted by a semi-structured interview guide, residents talked about their connection to Amherst, the issues that matter to them, and the ways they participate in the community. They also completed graphic elicitation activities where they identified important issues, their advocacy network, and a timeline of their civic identity development.  Thematic analyses, which are ongoing, involved all members of the research team independently reviewing the interviews and meeting to group similar responses into shared themes.
    }
    {
    Our early findings show four common themes in how residents talked about civic engagement. The I/nstitutional theme describes how government, money, and other larger systems affect residents{\textquoteright} ability to participate. The Inter/generational theme shows how family, education, and experiences across generations shape community involvement. The Team/work theme focuses on the importance of working with others, as well as the challenges that can make teamwork difficult. Finally, the Co/mpound theme describes how residents combine different approaches, such as voting, building relationships, sharing information, and joining community groups, to make a difference. Across these themes, residents often felt they had the least influence over the issues they considered most important.
    }
    {
    These early findings show that civic engagement is about more than voting or attending town meetings. How people participate is also connected to their relationships, life experiences, access to resources, and belief that they can make a difference. As we continue our analysis, we will compare the three groups to better understand how the Civic Academy may shape residents{\textquoteright} connection to Amherst, confidence in their ability to create change, and views of local government. Through our ongoing research-practice partnership, these findings may also help guide future Civic Academy programs.
    }

    \newpage
    